% Generic exact correlators and spectral representation in SI units

沿用上一小节的实标量场 $\varphi(x)$ 和四次自相互作用。这样做的好处是，传播子、顶角和 LSZ 的记号都已经出现过；现在只改变问题本身：不再把相互作用截断在某个微扰阶，而是问完整相关函数必须满足哪些恒等式。$\varphi$ 的 SI 单位仍为 $[\varphi]=\sqrt{\mathrm{J/m}}$，静质量 $m_\varphi$ 的单位为 kg，无量纲耦合仍记为 $g_4$。相应拉格朗日密度是
\begin{equation}
 \boxed{
 \Lag[\varphi]
 =\frac12\partial_\mu\varphi\,\partial^\mu\varphi
 -\frac12\frac{m_\varphi^2c^2}{\hbar^2}\varphi^2
 -\frac{g_4}{4!\,\hbar c}\varphi^4.}
 \label{eq:SD-scalar-L-SI-r12}
\end{equation}
$\Lag$ 的单位为 $\mathrm{J\,m^{-3}}$。不同阶相关函数可以统一由外源泛函微分产生。定义带源作用量
\begin{equation*}
 S[\varphi]=\frac1c\int\dd^4 x\,\Lag[\varphi].
\end{equation*}
再加入一个复数值函数 $\mathcal J(x)$，只把它当作产生场插入的数学记账变量，而不把它解释成实验电流。源作用量和归一化生成泛函分别定义为
\begin{align}
 S_{\rm src}[\varphi,\mathcal J]
 &=\frac1c\int\dd^4 x\,\mathcal J(x)\varphi(x),
 \label{eq:SD-source-action-dp76}\\
 Z[\mathcal J]
 &=\frac1{\mathcal N}\int\mathcal D\varphi\,
 \exp\!\left\{\frac{\ii}{\hbar}
 [S[\varphi]+S_{\rm src}[\varphi,\mathcal J]]\right\},
 \qquad Z[0]=1.
 \label{eq:SD-generating-functional-dp76}
\end{align}
源的量纲由 $\mathcal J\varphi$ 必须具有能量密度量纲固定为
\begin{equation*}
 [\mathcal J]=\mathrm{J^{1/2}m^{-5/2}}.
\end{equation*}
对 $Z[\mathcal J]$ 作一次泛函微分会在积分中带下一次 $\varphi(x)$ 插入，因此 $n$ 点时间有序相关函数可统一记为
\begin{equation}
 G_n(x_1,\ldots,x_n)
 \equiv
 \left.
 \left(\frac{\hbar c}{\ii}\right)^n
 \frac{\delta^nZ}
 {\delta\mathcal J(x_1)\cdots\delta\mathcal J(x_n)}
 \right|_{\mathcal J=0}.
 \label{eq:SD-correlator-source-definition-dp76}
\end{equation}
这里的 $G_n$ 是同一个量子理论在不同插入点数下的相关函数，不是彼此独立定义的动力学对象。

把经典 Euler--Lagrange 算符定义为
\begin{equation}
 \boxed{
 \mathfrak E[\varphi](x)
 \equiv c\frac{\delta S}{\delta\varphi(x)}
 =-\left(\Box+\frac{m_\varphi^2c^2}{\hbar^2}\right)\varphi(x)
 -\frac{g_4}{3!\,\hbar c}\varphi^3(x).}
 \label{eq:SD-euler-operator-SI-r16}
\end{equation}
归一化生成泛函记为 $Z[\mathcal J]$，并取 $Z[0]=1$。场空间测度在平移下不变这一条件，把“经典运动方程”提升为作用在整个生成泛函上的精确函数微分方程：
\begin{equation}
 \boxed{
 \left[
 \mathfrak E\!\left[
 \frac{\hbar c}{\ii}\frac{\delta}{\delta\mathcal J}
 \right](x)
 +\mathcal J(x)
 \right]Z[\mathcal J]=0.}
 \label{eq:SD-master-SI-r12}
\end{equation}
这就是 Schwinger--Dyson 母方程。它的物理内容是：经典场方程中的每个 $\varphi$ 因子，在量子理论里都变成对外源的泛函微分，所以一个固定点数的相关函数会与更高点数相关函数耦合。对母方程再作一次源微分并令 $\mathcal J=0$，最先得到
\begin{equation}
 \boxed{
 \left(\Box_x+\frac{m_\varphi^2c^2}{\hbar^2}\right)G_2(x,y)
 +\frac{g_4}{3!\,\hbar c}G_4(x,x,x,y)
 =-\ii\hbar c\,\delta^{(4)}(x-y).}
 \label{eq:SD-two-four-hierarchy-dp76}
\end{equation}
因此相互作用理论的二点函数通常不能独立闭合；继续微分会依次得到 $G_4$、$G_6$ 等更高相关函数的方程，形成无限层级。

对二点函数本身，另一种精确组织方式是把所有单粒子不可约二点插入重求和。记自由约化传播子为 $G_0(p)$、完整约化传播子为 $G(p)$。定义标量自能 $\Sigma_\varphi(p)$ 为所有不能通过切断一条内部标量线而分离的二点单粒子不可约插入之和。由于 $[G_0]=[p]^{-2}$，有 $[\Sigma_\varphi]=[p]^2$。按 1PI 自能插入次数组织完整二点函数，先得到形式级数
\begin{equation}
 \boxed{
 G
 =G_0+G_0\Sigma_\varphi G_0
 +G_0\Sigma_\varphi G_0\Sigma_\varphi G_0+\cdots .}
 \label{eq:exact-dyson-series-scalar-r68}
\end{equation}
把这一级数重求和后，等价的逆传播子形式为
\begin{equation}
 \boxed{G^{-1}(p)=G_0^{-1}(p)-\Sigma_\varphi(p).}
 \label{eq:exact-dyson-scalar-SI-r15}
\end{equation}
Dirac 场和规范场的对应关系分别提升为自旋矩阵和 Lorentz 张量逆核；标量情形给出这一逆核结构的最小代数原型。

Dyson 方程告诉我们“相互作用怎样改写逆传播子”，但还没有说明完整传播子允许包含哪些谱成分。为此取精确真空 $|\Omega\rangle$，并令 $|n\rangle$ 是完整 Hamilton 算符和四动量算符的共同谱态。对厄米标量场定义 Wightman 二点函数
\begin{equation*}
 W(x)\equiv\langle\Omega|\hat\varphi(x)\hat\varphi(0)|\Omega\rangle.
\end{equation*}
在两个场之间插入完整态的单位分解，并使用平移协变性，得到离散记号下的精确谱和
\begin{equation}
 \boxed{
 W(x)
 =\sum_n \ee^{-\ii p_n\cdot x/\hbar}
 |\langle\Omega|\hat\varphi(0)|n\rangle|^2.}
 \label{eq:W-spectral-sum-SI-r12}
\end{equation}
连续谱时求和符号同时代表相应 Lorentz 不变积分。矩阵元以模方出现，所以谱权重非负。把具有同一不变量
\begin{equation*}
 \zeta\equiv p_n^2,\qquad [\zeta]=[p]^2,
\end{equation*}
的态合并，定义 Källén--Lehmann 谱密度 $\rho_{\rm KL}(\zeta)$。采用全书的物理四动量归一化，Wightman 函数可写成
\begin{equation}
 \boxed{
 W(x)=\int_0^\infty\dd\zeta\,\rho_{\rm KL}(\zeta)
 \int\frac{\dd^3 p}{(2\pi\hbar)^3}\frac{\hbar^2 c}{2p^0}
 \ee^{-\ii p\cdot x/\hbar},
 \qquad p^0=\sqrt{\mathbf p^2+\zeta}.}
 \label{eq:KL-Wightman-mass-shell-r15}
\end{equation}
这里 $p^0$ 具有动量量纲，而不是能量量纲。谱密度满足
\begin{equation}
 \boxed{
 \rho_{\rm KL}(\zeta)\ge0,
 \qquad
 \int_0^\infty\dd\zeta\,\rho_{\rm KL}(\zeta)=1,
 \qquad [\rho_{\rm KL}]=[p]^{-2}.}
 \label{eq:KL-positive-normalization-r16}
\end{equation}
第一条来自完整态矩阵元的模方，第二条由等时正则对易关系固定。为清楚展示归一化来源，正文先记录两种计算同一等时导数的目标式：谱表示给
\begin{align}
 \left.\partial_{x^0}C(x-y)\right|_{x^0=y^0}
 &=-\ii\hbar c
 \left[\int_0^\infty\dd\zeta\,\rho_{\rm KL}(\zeta)\right]
 \delta^{(3)}(\mathbf x-\mathbf y),
 \label{eq:KL-commutator-spectral-derivative-r15}
\end{align}
而 $C(x-y)\equiv\langle\Omega|[\hat\varphi(x),\hat\varphi(y)]|\Omega\rangle$ 与正则对易关系给
\begin{equation}
 \left.\partial_{x^0}C(x-y)\right|_{x^0=y^0}
 =-\ii\hbar c\,\delta^{(3)}(\mathbf x-\mathbf y).
 \label{eq:KL-commutator-canonical-r15}
\end{equation}
比较两式即固定\eqnrefs{eq:KL-positive-normalization-r16} 的总谱权重；详细质量壳代数放在节末推导区。

时间排序把正、负时间的 Wightman 函数拼成 Feynman 二点函数。其约化动量空间形式是
\begin{equation}
 \boxed{
 G_{F,\rm red}(p)
 =\int_0^\infty\dd\zeta\,
 \rho_{\rm KL}(\zeta)
 \frac{\ii}{p^2-\zeta+\ii0^+}.}
 \label{eq:KL-SI-r12}
\end{equation}
因此每一个精确中间态质量都对应一个自由型 Feynman 分母，而相互作用的信息全部进入正谱密度。若场与稳定单粒子态有非零重叠，谱密度包含孤立 $\delta$ 峰：
\begin{equation}
 \boxed{
 \rho_{\rm KL}(\zeta)
 =Z\,\delta(\zeta-m_{\rm phys}^2c^2)
 +\rho_{\rm KL,cont}(\zeta).}
 \label{eq:rho-pole-cont-SI-r12}
\end{equation}
其中 $Z\ge0$ 是场对该稳定单粒子态的谱权重，$m_{\rm phys}$ 是物理极点质量，$\rho_{\rm KL,cont}$ 是从最低多粒子阈值开始的连续谱。代回谱表示得到
\begin{equation}
 \boxed{
 G_{F,\rm red}(p)
 =\frac{\ii Z}{p^2-m_{\rm phys}^2c^2+\ii0^+}
 +\int_{\zeta_{\rm th}}^\infty\dd\zeta\,
 \rho_{\rm KL,cont}(\zeta)
 \frac{\ii}{p^2-\zeta+\ii0^+}.}
 \label{eq:KL-pole-plus-cont-SI-r12}
\end{equation}
若最轻连续谱通道由两个质量为 $m$ 的粒子组成，则 $\zeta_{\rm th}=4m^2c^2$。这个阈值并非只是一条抽象的“支割起点”：四维二体相空间本身已经带有无量纲开阈因子
\begin{equation}
 \boxed{
 \beta_2(\zeta)
 =\sqrt{1-\frac{4m^2c^2}{\zeta}}\,
 \Theta(\zeta-4m^2c^2).}
 \label{eq:two-particle-threshold-factor-dp68}
\end{equation}
具体的 $\rho_{\rm KL,cont}$ 还要乘上真空到该多粒子态的矩阵元模方和相应态密度，但无论相互作用细节怎样，二粒子通道都不能在 $\zeta<4m^2c^2$ 出现。\figref{fig:qft-two-particle-threshold-dp68} 直接画出\eqnrefs{eq:two-particle-threshold-factor-dp68} 的开阈行为：连续谱不是在阈值处突然获得有限常数权重，而是先由相空间从零开启。

\begin{figure}[htbp]
 \centering
 \includegraphics[width=0.72\textwidth]{qft_two_particle_threshold_dp68.pdf}
 \caption{两个等质量粒子连续谱的普适相空间开阈因子。横轴 $x=\zeta/(m^2c^2)$；虚线 $x=4$ 是二体阈值。实际 Källén--Lehmann 连续谱还包含相互作用矩阵元权重，因此本图只展示由运动学固定的阈值部分。}
 \label{fig:qft-two-particle-threshold-dp68}
\end{figure}

孤立极点与支割线因此不是两个额外模型，而是同一个精确二点函数的离散谱和连续谱贡献；LSZ 正是利用孤立极点把场插入约化为渐近单粒子态。

同一解析思想也控制线性响应。对任意厄米玻色可观测量 $\hat A$，推迟响应函数定义为
\begin{equation}
 \boxed{
 \chi_R(t)
 =-\frac{\ii}{\hbar}\Theta(t)
 \langle[\hat A(t),\hat A(0)]\rangle.}
 \label{eq:retarded-susceptibility-def-r68}
\end{equation}
因果支撑 $\chi_R(t<0)=0$ 使其 Fourier 变换在复频率上半平面解析。只要高频衰减足以使闭合轮廓积分成立，实轴边界值便满足
\begin{equation}
 \boxed{
 \begin{aligned}
 \operatorname{Re}\chi_R(\omega)
 &=\frac1\pi\mathcal P\int_{-\infty}^{\infty}\dd\omega'
 \frac{\operatorname{Im}\chi_R(\omega')}{\omega'-\omega},\\
 \operatorname{Im}\chi_R(\omega)
 &=-\frac1\pi\mathcal P\int_{-\infty}^{\infty}\dd\omega'
 \frac{\operatorname{Re}\chi_R(\omega')}{\omega'-\omega}.
 \end{aligned}}
 \label{eq:kramers-kronig-r68}
\end{equation}
这就是 Kramers--Kronig 关系：吸收与色散并不是两个独立经验函数，而是同一个因果响应的虚部和实部。

稳定粒子的极点位于物理支的实质量壳。可衰变激发与连续谱混合后，以能量 $E$ 为变量的推迟预解式在共振附近可先写成精确的局部分解
\begin{equation*}
 G_R(E)
 =\frac{Z_R}{E-E_R+\ii\Gamma_E/2}+G_{\rm reg}(E),
\end{equation*}
其中 $G_{\rm reg}$ 在所考察共振邻域解析。令 $\Delta_R$ 表示从 $E_R$ 到最近其他极点、阈值支点或背景显著变化尺度的能量距离，并定义
\begin{equation*}
 \epsilon_R
 =\frac{\max(|E-E_R|,\Gamma_E)}{\Delta_R},
 \qquad
 \epsilon_{\rm bg}
 =\frac{|E-E_R+\ii\Gamma_E/2|\,|G_{\rm reg}(E_R)|}{|Z_R|}.
\end{equation*}
当 $\epsilon_R\ll1$ 且 $\epsilon_{\rm bg}\ll1$ 时，解析背景的变化与幅度均可作为次阶修正，于是得到单极点近似
\begin{equation}
 \boxed{
 G_R(E)
 =\frac{Z_R}{E-E_R+\ii\Gamma_E/2}
 \left[1+O(\epsilon_R,\epsilon_{\rm bg})\right].}
 \label{eq:resonance-pole-SI-r12}
\end{equation}
$E_R$ 与 $\Gamma_E$ 均以 J 为单位；极点项在时域给出振幅 $\exp[-\Gamma_E t/(2\hbar)]$，因此概率衰减寿命为 $\tau=\hbar/\Gamma_E$。在相应黎曼支上解析延拓后，共振极点位于非物理支。纳米光子学准正常模、极化激元共振以及电子自能引起的能移与线宽都可用这一受控极点展开描述；当共振接近另一极点或连续谱阈值、使 $\epsilon_R$ 不再小，单极点形式必须恢复为多极点或极点--支割线的完整谱表示。


\begin{exbox}[由共振极点直接读出线宽与寿命]
在 $\epsilon_R,\epsilon_{\rm bg}\ll1$ 的单极点区间内，若探测链在线性区把响应强度乘以能量无关增益 $C_{\mathcal I}>0$，定义 $\mathcal I(E)\equiv C_{\mathcal I}|G_R(E)|^2$，则
\begin{equation}
 \mathcal I(E)
 =C_{\mathcal I}|G_R(E)|^2
 \simeq
 \frac{C_{\mathcal I}|Z_R|^2}{(E-E_R)^2+(\Gamma_E/2)^2}.
 \label{eq:resonance-lorentzian-example-dp48}
\end{equation}
峰值位于 $E=E_R$。要求信号降到峰值的一半，有
\begin{align*}
 (E-E_R)^2+(\Gamma_E/2)^2
 &=2(\Gamma_E/2)^2,\\
 |E-E_R|&=\Gamma_E/2.
\end{align*}
因此两个半高点相距
\begin{equation*}
 \Delta E_{\rm FWHM}=\Gamma_E.
\end{equation*}
另一方面，同一极点在时域给出归一化概率
$P(t)=P(0)\ee^{-\Gamma_E t/\hbar}$，故寿命
\begin{equation*}
 \tau=\frac{\hbar}{\Gamma_E}.
\end{equation*}
这说明频域线宽与时域寿命不是两个独立经验量，而是同一个复极点的实频与时间表示。对窄带光学或电子--光子共振，还可定义无量纲品质因子 $Q_R\equiv E_R/\Gamma_E$；只有在单极点条件成立时，$Q_R$ 才能由这一 Lorentz 型线宽直接解释。
\end{exbox}

% -----------------------------------------------------------------------------
% -----------------------------------------------------------------------------

\begin{exbox}[Dyson 自能与寿命]
精确二点函数的谱表示告诉我们，稳定单粒子对应实轴上的孤立极点；一旦该激发能向其它在壳通道衰变，自能获得虚部，极点离开实轴。相应的可计算对象是能量变量 $E$ 下的推迟 Green 函数
\begin{equation*}
 G^R(E)=\frac{1}{E-E_0-\Sigma^R(E)}.
\end{equation*}
在一个窄共振附近，把自能在 $E=E_r$ 邻域展开成
\begin{equation*}
 \Sigma^R(E)
 \simeq \Delta E
 +(1-Z^{-1})(E-E_r)
 -\frac{\ii\Gamma}{2Z},
\end{equation*}
其中 $\Gamma>0$ 具有能量量纲，$Z$ 是极点留数。把实部吸收到重整化共振能量 $E_r=E_0+\Delta E$ 后，Green 函数约化为
\begin{equation}
 \boxed{
 G^R(E)\simeq
 \frac{Z}{E-E_r+\ii\Gamma/2}.}
 \label{eq:dyson-pole-breit-wigner-dp65}
\end{equation}
相应谱函数定义为 $A(E)=-2\operatorname{Im}G^R(E)$，因此
\begin{equation}
 \boxed{
 A(E)=
 \frac{Z\Gamma}{(E-E_r)^2+(\Gamma/2)^2}.}
 \label{eq:dyson-lorentzian-spectral-dp65}
\end{equation}
若该窄共振几乎耗尽单粒子谱权重，则 $Z\simeq1$，并且
\begin{equation*}
 \int_{-\infty}^{\infty}\frac{\dd E}{2\pi}A(E)=Z.
\end{equation*}
所以 $\Gamma$ 不是随意画在谱线上的“展宽参数”，而是由同一个复极点的虚部固定：$\Gamma=-2\,\operatorname{Im}E_\star$，同时也是该 Lorentzian 的半峰全宽。

再把\eqnrefs{eq:dyson-pole-breit-wigner-dp65} Fourier 变回时间。对 $t>0$ 闭合下半平面，只拾取
\begin{equation*}
 E_\star=E_r-\frac{\ii\Gamma}{2}
\end{equation*}
这个极点，得到
\begin{equation*}
 G^R(t)\simeq
 -\frac{\ii Z}{\hbar}\,\theta(t)
 \exp\!\left(-\frac{\ii E_rt}{\hbar}\right)
 \exp\!\left(-\frac{\Gamma t}{2\hbar}\right),
\end{equation*}
因此振幅寿命尺度是 $2\hbar/\Gamma$，而概率或能量按
\begin{equation*}
 \ee^{-\Gamma t/\hbar}
\end{equation*}
衰减，对应通常定义的寿命
\begin{equation}
 \boxed{\tau=\frac{\hbar}{\Gamma}.}
 \label{eq:width-lifetime-si-dp65}
\end{equation}

作为数量级练习，若一个光学共振位于 $E_r=1.55\,\mathrm{eV}$，品质因子 $Q=100$，并按 $Q=E_r/\Gamma$ 定义能量线宽，则
\begin{align*}
 \Gamma&=15.5\,\mathrm{meV},\\
 \tau&=\frac{\hbar}{\Gamma}
 \simeq42.5\,\mathrm{fs}.
\end{align*}
“复极点 $\leftrightarrow$ Lorentzian 频谱 $\leftrightarrow$ 指数时间衰减”是推迟预解式的一般解析结构；有损光学正常模、QNM Green 张量和电子能量损失谱都由相应自能的复极点实现同一结构。


\medskip

稳定粒子对应实轴上的孤立极点；当该态能够衰变时，推迟 Green 函数的极点移入下半复能量平面。\eqnrefs{eq:dyson-pole-breit-wigner-dp65} 已把窄共振附近的分母写成 Breit--Wigner 形式。对能量变量作 Fourier 变换时，下半平面的极点直接产生时间域指数衰减，并给出\eqnrefs{eq:width-lifetime-si-dp65} 的 $\tau=\hbar/\Gamma$。这一近似要求共振孤立、$\Gamma\ll E_r$，且自能在宽度 $\Gamma$ 的能量范围内变化缓慢。


从 Dyson 分母
\begin{equation*}
 D(E)=E-E_0-\Sigma^R(E)
\end{equation*}
开始。把自能实部在 $E_r$ 附近展开：
\begin{equation*}
 \operatorname{Re}\Sigma^R(E)
 \simeq \operatorname{Re}\Sigma^R(E_r)
 +(E-E_r)\,\partial_E\operatorname{Re}\Sigma^R(E_r).
\end{equation*}
令
\begin{equation*}
 E_r-E_0-\operatorname{Re}\Sigma^R(E_r)=0,
 \qquad
 Z^{-1}=1-\partial_E\operatorname{Re}\Sigma^R(E_r),
\end{equation*}
并定义
\begin{equation*}
 \Gamma=-2Z\operatorname{Im}\Sigma^R(E_r)>0.
\end{equation*}
于是
\begin{align*}
 D(E)
 &\simeq Z^{-1}(E-E_r)+\ii\frac{\Gamma}{2Z},\\
 G^R(E)=\frac1{D(E)}
 &\simeq\frac{Z}{E-E_r+\ii\Gamma/2},
\end{align*}
即\eqnrefs{eq:dyson-pole-breit-wigner-dp65}。

取虚部：
\begin{align*}
 \operatorname{Im}\frac{1}{x+\ii a}
 &=-\frac{a}{x^2+a^2},\\
 A(E)=-2\operatorname{Im}G^R(E)
 &=\frac{Z\Gamma}{(E-E_r)^2+(\Gamma/2)^2},
\end{align*}
得到\eqnrefs{eq:dyson-lorentzian-spectral-dp65}。再计算时间核
\begin{equation*}
 G^R(t)=\int_{-\infty}^{\infty}\frac{\dd E}{2\pi\hbar}
 \ee^{-\ii Et/\hbar}G^R(E).
\end{equation*}
对 $t>0$，因 $\ee^{-\ii Et/\hbar}$ 在下半平面衰减，闭合下半圆；唯一极点是 $E_\star=E_r-\ii\Gamma/2$。留数定理给
\begin{equation*}
 G^R(t)\simeq
 -\frac{\ii Z}{\hbar}\,
 \ee^{-\ii E_rt/\hbar}\ee^{-\Gamma t/(2\hbar)}.
\end{equation*}
振幅模方因而为
\begin{equation*}
 |G^R(t)|^2\simeq
 \frac{|Z|^2}{\hbar^2}\ee^{-\Gamma t/\hbar}.
\end{equation*}
若把寿命定义为概率下降到 $\ee^{-1}$ 的时间，便得到\eqnrefs{eq:width-lifetime-si-dp65}。
\end{exbox}




\begin{derivation}[从\eqnrefs{eq:SD-scalar-L-SI-r12}到\eqnrefs{eq:SD-master-SI-r12}]
由正文\eqnrefs{eq:SD-generating-functional-dp76}，在泛函积分中作无穷小平移 $\varphi\to\varphi+\delta\varphi$。测度平移不变且场空间边界项消失时，
\begin{equation*}
 0=
 \int\mathcal D\varphi\,
 \frac{\delta}{\delta\varphi(x)}
 \exp\!\left\{
 \frac{\ii}{\hbar}(S+S_{\rm src})
 \right\}.
\end{equation*}
把导数作用到指数，
\begin{equation*}
 0=
 \int\mathcal D\varphi\,
 \frac{\ii}{\hbar}
 \left[
 \frac{\delta S}{\delta\varphi(x)}
 +\frac1c\mathcal J(x)
 \right]
 \ee^{\ii(S+S_{\rm src})/\hbar}.
\end{equation*}
正文\eqnrefs{eq:SD-euler-operator-SI-r16} 给出
\begin{equation*}
 c\frac{\delta S}{\delta\varphi(x)}
 =\mathfrak E[\varphi](x).
\end{equation*}
另一方面，由正文的源生成泛函直接求导，
\begin{equation*}
 \frac{\delta Z}{\delta\mathcal J(x)}
 =\frac{\ii}{\hbar c}
 \langle\varphi(x)\rangle_{\mathcal J}Z,
\end{equation*}
所以积分中的一次场插入可由
\begin{equation*}
 \varphi(x)\longrightarrow
 \frac{\hbar c}{\ii}\frac{\delta}{\delta\mathcal J(x)}
\end{equation*}
替换。将同一替换用于 $\mathfrak E[\varphi]$ 中的各次幂，便有
\begin{equation*}
 \left[
 \mathfrak E\!\left[
 \frac{\hbar c}{\ii}\frac{\delta}{\delta\mathcal J}
 \right](x)
 +\mathcal J(x)
 \right]Z[\mathcal J]=0,
\end{equation*}
即正文\eqnrefs{eq:SD-master-SI-r12}。
\end{derivation}

\begin{derivation}[从\eqnrefs{eq:SD-master-SI-r12}到\eqnrefs{eq:SD-two-four-hierarchy-dp76}]
母方程中“把场换成源导数”很紧凑，但第一次接触 Schwinger--Dyson 层级时，最容易混淆的是三件事：每次源微分带来什么 $\hbar,c,\ii$ 因子，$\varphi^3(x)$ 为什么产生四点函数，以及右端的 $\delta^{(4)}(x-y)$ 从哪里出现。为把这三件事分开，定义场插入算符
\begin{equation*}
\widehat\Phi_x
\equiv\frac{\hbar c}{\ii}\frac{\delta}{\delta\mathcal J(x)}.
\end{equation*}
正文\eqnrefs{eq:SD-correlator-source-definition-dp76} 因而等价于
\begin{equation*}
G_n(x_1,\ldots,x_n)
=\left.
\widehat\Phi_{x_1}\cdots\widehat\Phi_{x_n}Z[\mathcal J]
\right|_{\mathcal J=0}.
\end{equation*}
特别地，$Z[0]=1$ 给
\begin{equation*}
\left.\widehat\Phi_x\widehat\Phi_yZ\right|_0=G_2(x,y),
\qquad
\left.\widehat\Phi_x^3\widehat\Phi_yZ\right|_0=G_4(x,x,x,y).
\end{equation*}
第二式没有额外的 $3!$：正文 Euler--Lagrange 算符里的 $1/3!$ 已经来自对 $\varphi^4/4!$ 求一次场导数；这里的三次 $\widehat\Phi_x$ 只是把三个位于同一点 $x$ 的场插入逐个生成出来。

把正文\eqnrefs{eq:SD-euler-operator-SI-r16} 代入母方程，先把算符写开：
\begin{equation*}
\left[
-\left(\Box_x+\frac{m_\varphi^2c^2}{\hbar^2}\right)\widehat\Phi_x
-\frac{g_4}{3!\,\hbar c}\widehat\Phi_x^3
+\mathcal J(x)
\right]Z[\mathcal J]=0.
\end{equation*}
现在从左边再作用一次 $\widehat\Phi_y$。前两个项中的源导数彼此可交换，所以
\begin{align*}
&-\left(\Box_x+\frac{m_\varphi^2c^2}{\hbar^2}\right)
\widehat\Phi_x\widehat\Phi_yZ
-\frac{g_4}{3!\,\hbar c}
\widehat\Phi_x^3\widehat\Phi_yZ\\
&\qquad
+\widehat\Phi_y\!\left[\mathcal J(x)Z\right]=0.
\end{align*}
真正产生接触项的是最后一个显式源因子。乘积法则给
\begin{align*}
\widehat\Phi_y\!\left[\mathcal J(x)Z\right]
&=\frac{\hbar c}{\ii}
\frac{\delta\mathcal J(x)}{\delta\mathcal J(y)}Z
+\mathcal J(x)\widehat\Phi_yZ\\
&=\frac{\hbar c}{\ii}\delta^{(4)}(x-y)Z
+\mathcal J(x)\widehat\Phi_yZ.
\end{align*}
令 $\mathcal J=0$ 后第二项消失，而 $Z[0]=1$，因此
\begin{equation*}
\left.\widehat\Phi_y[\mathcal J(x)Z]\right|_0
=-\ii\hbar c\,\delta^{(4)}(x-y).
\end{equation*}
同时利用上面的相关函数定义，母方程的一次源微分成为
\begin{equation*}
-\left(\Box_x+\frac{m_\varphi^2c^2}{\hbar^2}\right)G_2(x,y)
-\frac{g_4}{3!\,\hbar c}G_4(x,x,x,y)
-\ii\hbar c\,\delta^{(4)}(x-y)=0.
\end{equation*}
整体乘以 $-1$ 就得到正文\eqnrefs{eq:SD-two-four-hierarchy-dp76}：
\begin{equation*}
\left(\Box_x+\frac{m_\varphi^2c^2}{\hbar^2}\right)G_2(x,y)
+\frac{g_4}{3!\,\hbar c}G_4(x,x,x,y)
=-\ii\hbar c\,\delta^{(4)}(x-y).
\end{equation*}
因此层级不闭合的来源也已经显式可见：四次相互作用的经典运动方程含 $\varphi^3$，对二点函数方程而言它必然生成三个 $x$ 插入加一个 $y$ 插入的 $G_4$；接触 $\delta$ 则来自对外源本身的泛函微分，而不是相互作用顶角。
\end{derivation}

\begin{derivation}[从\eqnrefs{eq:exact-dyson-series-scalar-r68}到\eqnrefs{eq:exact-dyson-scalar-SI-r15}]
按自能插入次数排列完整标量二点函数。每个 $\Sigma_\varphi$ 两侧都由自由传播子连接，因此

\begin{align*}
 G
 &=G_0+G_0\Sigma_\varphi G_0
 +G_0\Sigma_\varphi G_0\Sigma_\varphi G_0+\cdots,\\
 &=G_0\left[
 1+\Sigma_\varphi G_0+(\Sigma_\varphi G_0)^2+\cdots
 \right].
\end{align*}
在几何级数作为形式微扰级数或相应逆算符存在的范围内，括号内可写成
\begin{equation*}
 \left(1-\Sigma_\varphi G_0\right)^{-1},
\end{equation*}
所以
\begin{equation*}
 G=G_0(1-\Sigma_\varphi G_0)^{-1}.
\end{equation*}
对标量动量空间核，$G_0(p)$ 与 $\Sigma_\varphi(p)$ 是同一固定四动量处的复数函数，因而彼此可交换。先取逆：
\begin{align*}
 G^{-1}
 &=\left[G_0(1-\Sigma_\varphi G_0)^{-1}\right]^{-1},\\
 &=(1-\Sigma_\varphi G_0)G_0^{-1},\\
 &=G_0^{-1}-\Sigma_\varphi G_0G_0^{-1},\\
 &=G_0^{-1}-\Sigma_\varphi.
\end{align*}
这就是\eqnrefs{eq:exact-dyson-scalar-SI-r15}。矩阵值或张量值传播子仍满足同一算符恒等式，但必须保留因子次序，不能使用标量交换性简化中间步骤。
\end{derivation}

\begin{derivation}[从\eqnrefs{eq:KL-Wightman-mass-shell-r15}到\eqnrefs{eq:KL-positive-normalization-r16}]
令 $r\equiv\mathbf x-\mathbf y$、
$\tau\equiv x^0-y^0$。对固定 $\zeta$，正文质量壳核为
\begin{equation*}
 W_\zeta(\tau,\mathbf r)
 =\int\frac{\dd^3p}{(2\pi\hbar)^3}
 \frac{\hbar^2c}{2p^0}
 \exp\!\left[-\frac{\ii}{\hbar}
 (p^0\tau-\mathbf p\cdot\mathbf r)\right],
 \qquad p^0=\sqrt{\mathbf p^2+\zeta}.
\end{equation*}
场对易函数中这一质量壳的贡献是
\[
 C_\zeta(\tau,\mathbf r)
 =W_\zeta(\tau,\mathbf r)-W_\zeta(-\tau,-\mathbf r).
\]
现在对 $x^0$ 求导；因为 $\partial_{x^0}\tau=1$，第一项给
\begin{align*}
 \left.\partial_{x^0}W_\zeta(\tau,\mathbf r)\right|_{\tau=0}
 &=-\frac{\ii\hbar c}{2}
 \int\frac{\dd^3p}{(2\pi\hbar)^3}
 \ee^{+\ii\mathbf p\cdot\mathbf r/\hbar}.
\end{align*}
第二项要同时注意前面的减号和 $-\tau$ 的链式法则：
\begin{align*}
 \left.\partial_{x^0}[-W_\zeta(-\tau,-\mathbf r)]\right|_{\tau=0}
 &=-\frac{\ii\hbar c}{2}
 \int\frac{\dd^3p}{(2\pi\hbar)^3}
 \ee^{-\ii\mathbf p\cdot\mathbf r/\hbar}.
\end{align*}
在第二个积分中令 $\mathbf p\mapsto-\mathbf p$，两项相同，因此
\begin{align*}
 \left.\partial_{x^0}C_\zeta(\tau,\mathbf r)\right|_{\tau=0}
 &=-\ii\hbar c
 \int\frac{\dd^3p}{(2\pi\hbar)^3}
 \ee^{+\ii\mathbf p\cdot\mathbf r/\hbar}\\
 &=-\ii\hbar c\,\delta^{(3)}(\mathbf r).
\end{align*}
注意 $p^0$ 已在求导时与核中的 $1/(2p^0)$ 精确约掉，所以这个等时结果与 $\zeta$ 无关。把所有质量壳按谱密度叠加：
\begin{align*}
 \left.\partial_{x^0}C(x-y)\right|_{x^0=y^0}
 =-\ii\hbar c\,
 \left[\int_0^\infty\dd\zeta\,\rho_{\rm KL}(\zeta)\right]
 \delta^{(3)}(\mathbf x-\mathbf y).
\end{align*}
另一方面，$\hat\pi=c^{-1}\partial_{x^0}\hat\varphi$，所以等时正则关系
$[\hat\varphi(\mathbf x),\hat\pi(\mathbf y)]=\ii\hbar\delta^{(3)}(\mathbf x-\mathbf y)$
给出同一个导数
\[
 \left.\partial_{x^0}C(x-y)\right|_{x^0=y^0}
 =-\ii\hbar c\,\delta^{(3)}(\mathbf x-\mathbf y).
\]
比较 $\delta^{(3)}$ 的系数得到
\[
 \int_0^\infty\dd\zeta\,\rho_{\rm KL}(\zeta)=1.
\]
再结合完整态插入中
$|\langle\Omega|\hat\varphi(0)|n\rangle|^2\ge0$，便得到正文
\eqnrefs{eq:KL-positive-normalization-r16} 的非负性与总权重归一化。
\end{derivation}

\begin{derivation}[从\eqnrefs{eq:KL-Wightman-mass-shell-r15}到\eqnrefs{eq:KL-SI-r12}]
固定谱参数 $\zeta$，定义
\[
 \varepsilon_{\mathbf p,\zeta}
 \equiv\sqrt{\mathbf p^2+\zeta}>0.
\]
正文\eqnrefs{eq:KL-Wightman-mass-shell-r15} 中对应这一 $\zeta$ 的正频质量壳核含时间因子
$\exp[-\ii\varepsilon_{\mathbf p,\zeta}x^0/\hbar]$。时间排序要求
\begin{align*}
 \Delta_F^{(\zeta)}(x)
 ={}&\Theta(x^0)\Delta_+^{(\zeta)}(x)
 +\Theta(-x^0)\Delta_+^{(\zeta)}(-x).
\end{align*}
关键问题是：什么四维动量分母的逆 Fourier 变换恰好产生这两个时间支路？对固定 $\mathbf p$ 考察
\begin{equation*}
 \frac{\ii}{(p^0)^2-\varepsilon_{\mathbf p,\zeta}^2+\ii0^+}.
\end{equation*}
其分布意义下的部分分式为
\begin{equation*}
 \frac{\ii}{(p^0)^2-\varepsilon^2+\ii0^+}
 =\frac{\ii}{2\varepsilon}
 \left[
 \frac1{p^0-\varepsilon+\ii0^+}
 -\frac1{p^0+\varepsilon-\ii0^+}
 \right].
\end{equation*}
因此正能极点位于 $p^0=+\varepsilon-\ii0^+$，负能极点位于
$p^0=-\varepsilon+\ii0^+$。逆 Fourier 因子是
$\exp[-\ii p^0x^0/\hbar]$：
当 $x^0>0$ 时，为使大半圆上的指数衰减，轮廓闭合到下半平面，只包围正能极点；当 $x^0<0$ 时闭合到上半平面，只包围负能极点。两次留数分别产生
\begin{align*}
 x^0>0:&\quad
 \frac{1}{2\varepsilon}
 \ee^{-\ii\varepsilon x^0/\hbar},\\
 x^0<0:&\quad
 \frac{1}{2\varepsilon}
 \ee^{+\ii\varepsilon x^0/\hbar},
\end{align*}
正好是 $\Delta_+^{(\zeta)}(x)$ 与
$\Delta_+^{(\zeta)}(-x)$ 的时间部分。把空间 Fourier 因子恢复后，固定 $\zeta$ 的时间排序核因此具有约化动量表示
\[
 \Delta_{F,\rm red}^{(\zeta)}(p)
 =\frac{\ii}{p^2-\zeta+\ii0^+}.
\]

K\"all\'en--Lehmann 表示只是把所有允许的不变量 $\zeta$ 按正谱权重线性叠加，所以
\begin{align*}
 G_{F,\rm red}(p)
 &=\int_0^\infty\dd\zeta\,
 \rho_{\rm KL}(\zeta)\Delta_{F,\rm red}^{(\zeta)}(p)\\
 &=\int_0^\infty\dd\zeta\,
 \rho_{\rm KL}(\zeta)
 \frac{\ii}{p^2-\zeta+\ii0^+},
\end{align*}
即正文\eqnrefs{eq:KL-SI-r12}。这里 $+\ii0^+$ 的作用已经由轮廓选择明确：它把正能极点放到实轴下方、负能极点放到实轴上方，从而编码“正频向未来传播、负频向过去传播”的时间排序边界条件。
\end{derivation}
\begin{derivation}[从\eqnrefs{eq:retarded-susceptibility-def-r68}到\eqnrefs{eq:kramers-kronig-r68}]
由正文\eqnrefs{eq:retarded-susceptibility-def-r68} 的因果支撑 $\chi_R(t)=0$（$t<0$），其 Fourier 变换

\begin{equation*}
 \chi_R(\omega)=\int_{-\infty}^{\infty}\dd t\,\ee^{+\ii\omega t}\chi_R(t)
\end{equation*}
在复频率上半平面解析：若 $\operatorname{Im}\omega>0$，因子 $\ee^{\ii\omega t}=\ee^{\ii\operatorname{Re}\omega t}\ee^{-\operatorname{Im}\omega t}$ 保证 $t>0$ 积分收敛。对上半平面闭合积分轮廓应用 Cauchy formula，取实轴边界值后得到
\begin{equation*}
 \chi_R(\omega)
 =\frac1{\pi\ii}
 \mathcal P\int_{-\infty}^{\infty}\dd\omega'
 \frac{\chi_R(\omega')}{\omega'-\omega}.
\end{equation*}
分别取实部与虚部：
\begin{align*}
 {
 \operatorname{Re}\chi_R(\omega)
 =\frac1\pi\mathcal P\int_{-\infty}^{\infty}\dd\omega'
 \frac{\operatorname{Im}\chi_R(\omega')}{\omega'-\omega}},\\
 {
 \operatorname{Im}\chi_R(\omega)
 =-\frac1\pi\mathcal P\int_{-\infty}^{\infty}\dd\omega'
 \frac{\operatorname{Re}\chi_R(\omega')}{\omega'-\omega}}.
 \end{align*}
两式正是正文\eqnrefs{eq:kramers-kronig-r68}。
\end{derivation}
